树级三部分子分布在软、共线边界发散，这说明“恰好观测到固定数目的部分子”不是红外安全的实验问题。有限能量阈值与有限角分辨率使一部分实胶子辐射无法与两部分子末态区分，而同一阶虚胶子修正又与这类未分辨实辐射共享完全相同的红外奇异结构。有限的可观测量因此必须把这些实验上简并的末态一起求和。

最简单的包容量是强子产生比
\begin{equation*}
 R(s)\equiv
 \frac{\sigma(e^+e^-\to\text{强子};s)}
 {\sigma(e^+e^-\to\mu^+\mu^-;s)},
 \qquad
 s=(p_{e^-}+p_{e^+})^2,
\end{equation*}
其中 $s$ 是总四动量平方，质心能量为 $E_{\rm cm}=c\sqrt s$。远离重夸克阈值和窄共振时，对每个活跃味定义
\begin{equation*}
 \epsilon_{m,q}=\frac{m_q^2c^2}{s}\ll1.
\end{equation*}
零质量近似保留 $O(\epsilon_{m,q}^0)$ 项。若还采用纯光子交换的简化玻恩权重，则需处在远低于 $Z$ 极点的区域，满足
\begin{equation*}
 \epsilon_Z\equiv\frac{s}{m_Z^2c^2}\ll1;
\end{equation*}
当这一条件不成立时，应把玻恩电流替换为完整 $\gamma/Z$ 电弱电流，而 QCD 实--虚红外抵消结构保持不变。在这两个受控简化下，单个夸克味的玻恩截面为
\begin{equation}
 {
 \sigma_{0,q}(s)
 =N_cQ_q^2\frac{4\pi\alpha^2\hbar^2}{3s}.}
 \label{eq:qcd-r-born-sigma-si-dp12}
\end{equation}
$Q_q$ 是以正基本电荷 $e$ 为单位的夸克电荷。由于 $s$ 具有动量平方量纲，$\hbar^2/s$ 因而具有面积量纲；这一量纲关系与物理四动量的 SI 定义一致。

在 $\order(\alpha_s)$，同一个包容截面包含两部分：两体末态 $e^+e^-\to q\bar q$ 的重整化虚修正，以及三体末态 $e^+e^-\to q\bar qg$ 的实辐射。维数正则化把两者都定义为 $\epsilon$ 的亚纯函数。紫外反项先消去虚振幅中的紫外极点，但虚修正和实辐射仍分别含有红外双极点与单极点；只有把两者按同一实验定义相加，这些红外极点才逐项抵消。

\begin{figure}[htbp]
\centering
\begin{tikzfig}[wfA]
  \begin{scope}[xshift=-4.2cm]
    \draw[fermion] (-2,1) -- (-.8,.3); \draw[fermion] (-.8,-.3) -- (-2,-1);
    \draw[photon] (-.8,.3) -- (.25,0); \draw[fermion] (.25,0) -- (2,1); \draw[fermion] (2,-1) -- (.25,0);
    \draw[decorate,decoration={coil,aspect=.45,segment length=4.8pt,amplitude=1.8pt}] (.85,.55) arc[start angle=130,end angle=-130,radius=.75];
    \node at (0,-1.55){virtual correction};
  \end{scope}
  \begin{scope}[xshift=3.5cm]
    \draw[fermion] (-2,1) -- (-.8,.3); \draw[fermion] (-.8,-.3) -- (-2,-1);
    \draw[photon] (-.8,.3) -- (.15,0); \draw[fermion] (.15,0) -- (2,1); \draw[fermion] (2,-1) -- (.15,0);
    \draw[decorate,decoration={coil,aspect=.45,segment length=4.8pt,amplitude=1.8pt}] (1.0,.55) -- (2.15,0);
    \node at (0,-1.55){real emission};
  \end{scope}
\end{tikzfig}
\caption{次领阶包容强子产生率必须同时包含虚修正和实验上未分辨的实胶子辐射。两部分分别含红外极点，而同一包容测量中的和是有限的。}
\label{fig:qcd-r-nlo-real-virtual-dp12}
\end{figure}

对每个近似无质量的活跃夸克味，实--虚抵消后的结果为
\begin{equation*}
 \sigma_q^{\rm NLO}
 =\sigma_{0,q}
 \left[1+\frac{\alpha_s(\mu_{\rm DR})}{\pi}
 +\order(\alpha_s^2)\right].
\end{equation*}
对所有活跃味求和，得到
\begin{equation}
 {
 R(s)=N_c\sum_{q\,\text{活跃}}Q_q^2
 \left[1+\frac{\alpha_s(\mu_{\rm DR})}{\pi}
 +\order(\alpha_s^2)\right].}
 \label{eq:qcd-r-nlo-final-dp12}
\end{equation}
这条公式把包容性、红外安全性与强耦合测量直接连接起来：最低阶给出颜色重数和电荷平方和，次领阶有限修正则对 $\alpha_s$ 敏感。由于 $R(s)$ 对强子末态做完全包容求和，它也可以通过前向电磁流关联函数的吸收部分来表示；这正是光学定理把“所有强子末态的和”改写成“一个前向关联函数的虚部”的一般机制。

有限结果依赖四个彼此不同的理论步骤。紫外重整化先定义有限的硬振幅；实辐射和虚修正单独仍然红外发散；实验分辨率要求对简并末态作包容求和；幺正性与 Kinoshita--Lee--Nauenberg 抵消最终使正则化参数从可观测量中消失。在这一阶，显式的 $\ln(\mu_{\rm DR}^2/s)$ 项相消；更高阶的尺度依赖通过运行的 $\alpha_s$ 和高阶系数函数重新出现，因此改变 $\mu_{\rm DR}$ 可以用来估计截断微扰级数的剩余尺度敏感性，但不能把它解释成新的物理参数。

适用域由三个相互独立的展开控制：夸克质量要求 $\epsilon_{m,q}\ll1$，固定阶强相互作用要求 $\alpha_s(\sqrt{s})\ll1$，而纯光子玻恩权重另要求 $\epsilon_Z\ll1$。有限质量修正从 $O(\epsilon_{m,q})$ 开始；弱中性流修正从 $O(\epsilon_Z)$ 进入。接近重夸克产生阈值时，速度和质量效应不再是小修正；接近 $Z$ 极点时必须恢复完整电弱传播子；接近窄强子共振时，局域固定阶部分子展开也不能直接描述谱峰。

